%!TEX TS-program = xelatex
%!TEX encoding = UTF-8 Unicode

%===============================================================================
% 模板配置
%   切换模板：将 job 改为 academic 即可使用学术版
%===============================================================================
\documentclass[11pt, a4paper, job]{dualcv}

\geometry{left=1.4cm, top=.8cm, right=1.4cm, bottom=1.8cm, footskip=.5cm}
\setbool{acvSectionColorHighlight}{true}



%===============================================================================
% 个人信息（共享）
%   求职版/学术版分别显示不同的 position 和链接，由 cls 自动处理
%===============================================================================
\photo[circle,noedge,left]{avatar}
\name{你的姓名}
\mobile{138-0000-0000}
\email{your-email@example.com}
% \homepage{www.example.com}
% \github{your-github-id}

\ifjob{\position{目标岗位名称}}
\ifacademic{
  % \position{计算机科学硕士申请}
  % \orcid{0000-0000-0000-0000}
  % \googlescholar{scholar-id}{你的姓名}
  \researchinterests{Research Interest A, Research Interest B}
}


%===============================================================================
\begin{document}
\makecvheader[R]
%===============================================================================


%###############################################################################
%
%                         共享内容（两版都会显示）
%
%###############################################################################

%-------------------------------------------------------------------------------
% 教育背景
%   \cveducation{学校}{学院/专业}{时间}{描述}
%-------------------------------------------------------------------------------
\cvsection{教育背景}

\begin{cventries}
  \cveducation
    {XX 大学}
    {XX 学院}
    {20XX.09 -- 20XX.06}
    {
      \begin{cvitems}
        \item 专业：XX 工程
        \item GPA：\textbf{X.X / 4.0}（专业前 X\%）
        \ifjob{\item 核心课程：数据结构、算法设计、操作系统、计算机网络、数据库系统}
        \ifacademic{\item 核心课程：数据结构、算法设计、操作系统、机器学习、计算机视觉、高等数学、线性代数、概率论}
      \end{cvitems}
    }
\end{cventries}


%###############################################################################
%
%                         求职版专属内容（仅 job 模式显示）
%
%###############################################################################
\ifjob{

%-------------------------------------------------------------------------------
% 实习经历
%   \cvinternship{公司}{职位}{时间}{描述}
%-------------------------------------------------------------------------------
\cvsection{实习经历}

\begin{cventries}
  \cvinternship
    {XX 公司}
    {XX 实习生}
    {20XX.07 -- 20XX.09}
    {
      \begin{cvitems}
        \item 参与 XX 系统后端服务开发，负责数据管道优化，将数据延迟从 Xmin 降至 Xmin
        \item 基于 XX 重构实时特征计算模块，支持每秒 X 万+ 事件处理
        \item 编写单元测试与集成测试，代码覆盖率达 XX\% 以上
      \end{cvitems}
    }
\end{cventries}

%-------------------------------------------------------------------------------
% 项目经历
%   \cvproject{项目名}{角色}{时间}{描述}
%-------------------------------------------------------------------------------
\cvsection{项目经历}

\begin{cventries}
  \cvproject
    {项目名称 A}
    {个人项目}
    {20XX.03 -- 20XX.08}
    {
      \begin{cvitems}
        \item 使用 XX 语言实现了一个支持 XX 的分布式系统
        \item 基于 XX 算法实现了核心功能，具备高可用性
        \item 采用 XX 进行节点间通信，HTTP API 对外提供服务，单节点 QPS 达 X000+
        \item 编写了完整的技术文档，GitHub Star XX+
      \end{cvitems}
    }

  \cvproject
    {项目名称 B}
    {后端开发}
    {20XX.06 -- 20XX.09}
    {
      \begin{cvitems}
        \item 独立负责后端设计与开发，使用 XX + XX 构建 RESTful API
        \item 实现 XX 认证，支持核心业务功能
        \item 引入 XX 缓存热门数据，接口响应时间从 XXms 降至 XXms
        \item 项目已上线，累计注册用户 XX+
      \end{cvitems}
    }

  \cvproject
    {项目名称 C}
    {课程项目}
    {20XX.01 -- 20XX.04}
    {
      \begin{cvitems}
        \item 使用 XX 实现 XX 模型，在 XX 数据集上准确率达 XX\%
        \item 利用数据增强有效缓解过拟合
        \item 使用 XX 进行模型推理加速，推理速度提升 X 倍
      \end{cvitems}
    }
\end{cventries}

%-------------------------------------------------------------------------------
% 技能证书
%   \cvskill{类别}{内容}
%-------------------------------------------------------------------------------
\cvsection{技能证书}

\begin{cvskills}
  \cvskill{编程语言}{Python / Go / Java / C++，熟悉常用数据结构和算法}
  \cvskill{后端框架}{Spring Boot / Gin / gRPC，了解 WebSocket、RESTful API 设计}
  \cvskill{数据库}{MySQL / PostgreSQL / Redis / MongoDB，熟悉索引优化与事务机制}
  \cvskill{中间件}{Kafka / RabbitMQ / Nginx / Docker / K8s，有实际项目使用经验}
  \cvskill{英语水平}{CET-X，能流畅阅读英文技术文档}
  \cvskill{资格证书}{XX 证书、XX 证书}
\end{cvskills}

%-------------------------------------------------------------------------------
% 荣誉奖励
%   \cvaward{年份}{获奖内容}
%-------------------------------------------------------------------------------
\cvsection{荣誉奖励}

\begin{cvhonors}
  \cvaward{20XX}{XX 大学 XX 奖学金}
  \cvaward{20XX}{XX 竞赛 XX 奖}
  \cvaward{20XX}{XX 竞赛 XX 奖}
\end{cvhonors}

}% === 求职版专属内容结束 ===


%###############################################################################
%
%                         学术版专属内容（仅 academic 模式显示）
%
%###############################################################################
\ifacademic{

%-------------------------------------------------------------------------------
% 研究经历
%   \cvresearchentry{课题名}{角色/导师}{时间}{描述}
%-------------------------------------------------------------------------------
\cvsection{研究经历}

\begin{cventries}
  \cvresearchentry
    {研究课题名称 A}
    {本科生科研助理 \textbar{} 指导教师：XX 教授}
    {20XX.03 -- 至今}
    {
      \begin{cvitems}
        \item 参与实验室 XX 课题，负责模型改进与实验设计
        \item 提出改进的 XX 结构，在 XX 数据集上指标提升 X.X\%
        \item 完成相关文献综述 XX+ 篇，撰写研究报告并在组会汇报
      \end{cvitems}
    }

  \cvresearchentry
    {研究课题名称 B}
    {课程研究项目}
    {20XX.01 -- 20XX.04}
    {
      \begin{cvitems}
        \item 使用 XX 实现 XX 模型，在 XX 数据集上准确率达 XX\%
        \item 利用数据增强有效缓解过拟合
        \item 使用 XX 进行模型推理加速，推理速度提升 X 倍
        \item 撰写技术报告，分析不同网络结构对性能的影响
      \end{cvitems}
    }
\end{cventries}

%-------------------------------------------------------------------------------
% 论文发表
%   \cvpublication{论文标题}{作者列表（本人 \textbf{加粗}）}{录用信息}
%-------------------------------------------------------------------------------
\cvsection{论文发表}

\begin{cvpublications}
  \cvpublication
    {Your Paper Title in English}
    {\textbf{Your Name}, Co-Author A, Co-Author B}
    {Journal Name, SCI X}

  \cvpublication
    {你的论文标题}
    {\textbf{你的姓名}, 合作者 A, 合作者 B}
    {已录用于 XX 会议/期刊，CCF-X}
\end{cvpublications}

%-------------------------------------------------------------------------------
% 竞赛奖项
%   \cvaward{年份}{获奖内容}
%-------------------------------------------------------------------------------
\cvsection{竞赛奖项}

\begin{cvhonors}
  \cvaward{20XX}{XX 竞赛 XX 奖}
  \cvaward{20XX}{XX 竞赛 XX 奖}
  \cvaward{20XX}{XX Workshop / Poster 展示}
\end{cvhonors}

%-------------------------------------------------------------------------------
% 技能与语言
%   \cvskill{类别}{内容}
%-------------------------------------------------------------------------------
\cvsection{技能与语言}

\begin{cvskills}
  \cvskill{编程语言}{Python / Go / Java / C++，熟悉常用数据结构和算法}
  \cvskill{科研工具}{PyTorch / TensorFlow / OpenCV / LaTeX / MATLAB}
  \cvskill{开发工具}{Git / Docker / Linux / MySQL / VS Code}
  \cvskill{英语水平}{CET-X，能流畅阅读英文技术文档}
\end{cvskills}

%-------------------------------------------------------------------------------
% 荣誉奖励
%   \cvaward{年份}{获奖内容}
%-------------------------------------------------------------------------------
\cvsection{荣誉奖励}

\begin{cvhonors}
  \cvaward{20XX}{XX 大学 XX 奖学金}
  \cvaward{20XX}{XX 奖学金}
\end{cvhonors}

}% === 学术版专属内容结束 ===


%===============================================================================
\end{document}
